\documentclass[10pt,letterpaper]{article}
\usepackage{cornell-notes}

\title{Annex D.26.02: Header Codecvt Synopsis}
\author{}
\date{}
\setDocTitle{Annex D.26.02: Header Codecvt Synopsis}
\setDocSubtitle{C++ 2024}
\setDocOwner{C++ 2024 Cornell Notes}
\setDocDate{\today}
\setCornellCollection{C++ 2024 Cornell Notes}
\setCornellUnitType{Annex}
\setCornellUnitNumber{D.26.02}
\setCornellUnitTitle{Header Codecvt Synopsis}
\hypersetup{pdftitle={Annex D.26.02: Header Codecvt Synopsis},pdfsubject={Cornell notes},pdfauthor={}}

\begin{document}
\maketitle
\begin{CornellOverviewBox}{Overview}
\textbf{Classification:} Normative annex\\[2pt]
\textbf{Learning objective:} Understand the role and technical guidance associated with Header <codecvt> synopsis.
\end{CornellOverviewBox}\begin{CornellSummaryBox}{Summary}
{\sffamily\bfseries\color{Accent} Summary}\par\vspace{3pt}Section D.26.2 focuses on Header <codecvt> synopsis. Use the listed grammar productions together with the semantic constraints and disambiguation rules referenced by the section. When applying it, preserve the distinction between mandatory requirements, recommendations, permissions, and behavior deliberately left undefined, unspecified, or implementation-defined.
\end{CornellSummaryBox}

\end{document}
